\textbf{Phần 1. Trắc nghiệm lựa chọn}
\setcounterpageref{ex}{0}

\begin{ex} %[5H8V48]
Một bể cá cảnh chứa được nhiều nhất $0{,}36\text{ m}^3$ nước. Hiện tại mực nước trong bể đạt $\dfrac{3}{4}$ chiều cao của bể. Người ta cần đổ thêm bao nhiêu lít nước nữa để bể đầy?
\choice
{$60\text{ lít}$}
{\True $90\text{ lít}$}
{$120\text{ lít}$}
{$270\text{ lít}$}
\loigiai{
Đổi: $0{,}36\text{ m}^3 = 360\text{ dm}^3 = 360\text{ lít}$.\\
Phần dung tích cần đổ thêm chiếm:
$$1 - \dfrac{3}{4} = \dfrac{1}{4}\text{ (bể)}.$$
Số lít nước cần đổ thêm là:
$$360 \times \dfrac{1}{4} = 90\text{ (lít)}.$$
\textbf{Đáp án của câu hỏi là B}
}
\end{ex}

\begin{ex} %[5H8V48]
Cho biểu thức $A = 2{,}4\text{ m}^3 - 650\text{ dm}^3 + 350\,000\text{ cm}^3$. Giá trị của $A$ đổi sang đơn vị đề-xi-mét khối là
\choice
{$1\,400\text{ dm}^3$}
{\True $2\,100\text{ dm}^3$}
{$1\,750\text{ dm}^3$}
{$2\,450\text{ dm}^3$}
\loigiai{
Đổi về cùng đơn vị đề-xi-mét khối:
- $2{,}4\text{ m}^3 = 2\,400\text{ dm}^3$.\\
- $350\,000\text{ cm}^3 = 350\text{ dm}^3$.\\
Giá trị của $A$ là:
$$2\,400 - 650 + 350 = 1\,750 + 350 = 2\,100\text{ (dm}^3\text{)}.$$
\textbf{Đáp án của câu hỏi là B}
}
\end{ex}

\begin{ex} %[5H8V48]
Một khối bê tông hình hộp chữ nhật có thể tích $1{,}8\text{ m}^3$. Biết diện tích một mặt đáy của khối bê tông là $2{,}4\text{ m}^2$. Chiều cao của khối bê tông đó là
\choice
{$0{,}65\text{ m}$}
{\True $0{,}75\text{ m}$}
{$0{,}8\text{ m}$}
{$0{,}72\text{ m}$}
\loigiai{
Chiều cao của khối bê tông là:
$$1{,}8 : 2{,}4 = 0{,}75\text{ (m)}.$$
\textbf{Đáp án của câu hỏi là B}
}
\end{ex}

\begin{ex} %[5H8V48]
Một bể chứa nước hình lập phương có cạnh $1{,}2\text{ m}$. Bể đang chứa $864\text{ lít}$ nước. Mực nước trong bể hiện cao bao nhiêu xăng-ti-mét?
\choice
{$50\text{ cm}$}
{\True $60\text{ cm}$}
{$70\text{ cm}$}
{$80\text{ cm}$}
\loigiai{
Đổi: $1{,}2\text{ m} = 12\text{ dm}$.\\
Diện tích đáy bể là:
$$12 \times 12 = 144\text{ (dm}^2\text{)}.$$
Đổi $864\text{ lít} = 864\text{ dm}^3$.\\
Chiều cao mực nước hiện tại là:
$$864 : 144 = 6\text{ (dm)}.$$
Đổi $6\text{ dm} = 60\text{ cm}$.\\
\textbf{Đáp án của câu hỏi là B}
}
\end{ex}

\begin{ex} %[5H8V48]
Một xe téc chứa $15\text{ m}^3$ dầu. Sau khi xả bớt dầu vào $40$ thùng phuy, mỗi thùng chứa $250\text{ lít}$, lượng dầu còn lại trong xe téc chiếm bao nhiêu phần trăm lượng dầu ban đầu?
\choice
{$25\%$}
{\True $33{,}33\%$}
{$50\%$}
{$66{,}67\%$}
\loigiai{
Đổi $15\text{ m}^3 = 15\,000\text{ dm}^3 = 15\,000\text{ lít}$.\\
Lượng dầu đã xả vào $40$ thùng phuy là:
$$250 \times 40 = 10\,000\text{ (lít)}.$$
Lượng dầu còn lại trong xe téc là:
$$15\,000 - 10\,000 = 5\,000\text{ (lít)}.$$
Tỉ số phần trăm lượng dầu còn lại so với ban đầu là:
$$5\,000 : 15\,000 \times 100\% \approx 33{,}33\%.$$
\textbf{Đáp án của câu hỏi là B}
}
\end{ex}

\begin{ex} %[5H8V48]
Xếp các khối lập phương nhỏ cạnh $2\text{ cm}$ để được một khối lập phương lớn có thể tích $0{,}008\text{ m}^3$. Hỏi cần dùng tất cả bao nhiêu khối lập phương nhỏ?
\choice
{$100$}
{\True $1\,000$}
{$10\,000$}
{$8\,000$}
\loigiai{
Thể tích của mỗi khối lập phương nhỏ là:
$$2 \times 2 \times 2 = 8\text{ (cm}^3\text{)}.$$
Đổi: $0{,}008\text{ m}^3 = 8\,000\text{ cm}^3$.\\
Số khối lập phương nhỏ cần dùng là:
$$8\,000 : 8 = 1\,000\text{ (khối)}.$$
\textbf{Đáp án của câu hỏi là B}
}
\end{ex}

\begin{ex} %[5H8V48]
Một căn phòng hình hộp chữ nhật có chiều dài $6\text{ m}$, chiều rộng $4\text{ m}$, chiều cao $3{,}5\text{ m}$. Biết trong phòng có $30$ học sinh và $1$ giáo viên. Biết mỗi người cần trung bình $2{,}5\text{ m}^3$ không khí để thở. Hỏi thể tích không khí trong phòng có đủ chuẩn cho cả lớp không và dư hay thiếu bao nhiêu mét khối?
\choice
{Thiếu $6{,}5\text{ m}^3$}
{\True Dư $6{,}5\text{ m}^3$}
{Thiếu $4\text{ m}^3$}
{Dư $9\text{ m}^3$}
\loigiai{
Thể tích không khí trong phòng là:
$$6 \times 4 \times 3{,}5 = 84\text{ (m}^3\text{)}.$$
Tổng số người trong phòng là:
$$30 + 1 = 31\text{ (người)}.$$
Lượng không khí cần thiết cho $31$ người là:
$$31 \times 2{,}5 = 77{,}5\text{ (m}^3\text{)}.$$
Vì $84 > 77{,}5$ nên không khí đủ và còn dư:
$$84 - 77{,}5 = 6{,}5\text{ (m}^3\text{)}.$$
\textbf{Đáp án của câu hỏi là B}
}
\end{ex}

\begin{ex} %[5H8V48]
Một vòi nước chảy vào một bể không có nước, mỗi phút chảy được $45\text{ lít}$. Cùng lúc đó, ở đáy bể có một lỗ rò rỉ làm thất thoát $15\text{ lít}$ mỗi phút. Sau bao lâu thì bể có dung tích $3{,}6\text{ m}^3$ sẽ đầy nước?
\choice
{$100\text{ phút}$}
{\True $120\text{ phút}$}
{$90\text{ phút}$}
{$150\text{ phút}$}
\loigiai{
Trong mỗi phút, lượng nước thực tế chảy vào bể là:
$$45 - 15 = 30\text{ (lít)}.$$
Đổi $3{,}6\text{ m}^3 = 3\,600\text{ dm}^3 = 3\,600\text{ lít}$.\\
Thời gian để bể đầy nước là:
$$3\,600 : 30 = 120\text{ (phút)}.$$
\textbf{Đáp án của câu hỏi là B}
}
\end{ex}

\begin{ex} %[5H8V48]
Một khối đá hình lập phương có cạnh dài $0{,}8\text{ m}$. Biết $1\text{ dm}^3$ đá đó cân nặng $2{,}6\text{ kg}$. Khối đá đó cân nặng bao nhiêu ki-lô-gam?
\choice
{$1\,331{,}2\text{ kg}$}
{\True $1\,331{,}2\text{ kg}$}
{$1\,664\text{ kg}$}
{$1\,064{,}96\text{ kg}$}
\loigiai{
Đổi $0{,}8\text{ m} = 8\text{ dm}$.\\
Thể tích của khối đá là:
$$8 \times 8 \times 8 = 512\text{ (dm}^3\text{)}.$$
Cân nặng của khối đá là:
$$512 \times 2{,}6 = 1\,331{,}2\text{ (kg)}.$$
\textbf{Đáp án của câu hỏi là B}
}
\end{ex}

\begin{ex} %[5H8V48]
Biết rằng $1\text{ lít}$ nước biển chứa khoảng $35\text{ g}$ muối. Hỏi trong một bể chứa $4\text{ m}^3$ nước biển có tất cả bao nhiêu ki-lô-gam muối?
\choice
{$14\text{ kg}$}
{\True $140\text{ kg}$}
{$1\,400\text{ kg}$}
{$70\text{ kg}$}
\loigiai{
Đổi: $4\text{ m}^3 = 4\,000\text{ dm}^3 = 4\,000\text{ lít}$.\\
Lượng muối có trong bể nước biển là:
$$35 \times 4\,000 = 140\,000\text{ (g)}.$$
Đổi $140\,000\text{ g} = 140\text{ kg}$.\\
\textbf{Đáp án của câu hỏi là B}
}
\end{ex}

\begin{ex} %[5H8V48]
Người ta thả một hòn non bộ vào một thùng nước hình hộp chữ nhật có diện tích đáy $30\text{ dm}^2$ đang chứa nước. Mực nước trong thùng dâng thêm $15\text{ cm}$. Thể tích của hòn non bộ đó là
\choice
{$450\text{ dm}^3$}
{\True $45\text{ dm}^3$}
{$4{,}5\text{ dm}^3$}
{$45\,000\text{ dm}^3$}
\loigiai{
Đổi: $15\text{ cm} = 1{,}5\text{ dm}$.\\
Thể tích nước dâng lên chính là thể tích của hòn non bộ:\\
$$V = 30 \times 1{,}5 = 45\text{ (dm}^3\text{)}.$$
\textbf{Đáp án của câu hỏi là B}
}
\end{ex}

\begin{ex} %[5H8V48]
Một bể nước dạng hình hộp chữ nhật có các kích thước trong lòng bể là: dài $2{,}5\text{ m}$, rộng $1{,}6\text{ m}$, cao $1{,}4\text{ m}$. Hiện tại bể đã có $2\,800\text{ lít}$ nước. Cần bơm thêm vào bể bao nhiêu mét khối nước nữa để đầy bể?
\choice
{$5{,}6\text{ m}^3$}
{\True $2{,}8\text{ m}^3$}
{$3{,}2\text{ m}^3$}
{$1{,}8\text{ m}^3$}
\loigiai{
Thể tích lòng bể là:
$$2{,}5 \times 1{,}6 \times 1{,}4 = 5{,}6\text{ (m}^3\text{)}.$$
Đổi $2\,800\text{ lít} = 2{,}8\text{ m}^3$.\\
Số mét khối nước cần bơm thêm là:
$$5{,}6 - 2{,}8 = 2{,}8\text{ (m}^3\text{)}.$$
\textbf{Đáp án của câu hỏi là B}
}
\end{ex}

\textbf{Phần 2. Trắc nghiệm đúng sai}
\setcounterpageref{ex}{0}

\begin{ex} %[5H8V48]
Một khối gỗ hình hộp chữ nhật có kích thước dài $1{,}2\text{ m}$, rộng $0{,}8\text{ m}$ và cao $0{,}5\text{ m}$. Người ta cắt khối gỗ thành các con xúc xắc hình lập phương cạnh $1\text{ dm}$. Xét tính đúng/sai của các nhận định sau:
\choiceTF
{\True Thể tích khối gỗ ban đầu là $0{,}48\text{ m}^3$}
{\True Thể tích khối gỗ ban đầu bằng $480\text{ dm}^3$}
{Mỗi con xúc xắc nhỏ có thể tích là $10\text{ dm}^3$}
{\True Nếu không hao hụt do mạch cưa thì cắt được đúng $480$ con xúc xắc}
\loigiai{
\begin{itemchoice}
\item Đúng. $V = 1{,}2 \times 0{,}8 \times 0{,}5 = 0{,}48\text{ m}^3$.
\item Đúng. $0{,}48\text{ m}^3 = 480\text{ dm}^3$.
\item Sai. Mỗi con xúc xắc có thể tích là $1 \times 1 \times 1 = 1\text{ dm}^3$.
\item Đúng. Số xúc xắc cắt được là $480 : 1 = 480$ (con xúc xắc). Các kích thước $12\text{ dm}, 8\text{ dm}, 5\text{ dm}$ đều chia hết cho $1\text{ dm}$.
\end{itemchoice}
}
\end{ex}

\begin{ex} %[5H8V48]
Một bể nước hình hộp chữ nhật chứa được tối đa $6\text{ m}^3$ nước. Buổi sáng người ta bơm vào $2\,400\text{ lít}$ nước. Buổi chiều dùng hết $\dfrac{1}{3}$ lượng nước hiện có trong bể. Xét tính đúng/sai của các nhận định sau:
\choiceTF
{\True Sau buổi sáng, lượng nước trong bể chiếm $40\%$ thể tích bể}
{\True Buổi chiều đã dùng hết $800\text{ lít}$ nước}
{Lượng nước còn lại sau buổi chiều là $1\,800\text{ lít}$}
{\True Muốn bể đầy nước cần bơm thêm $4\,400\text{ lít}$ nước nữa}
\loigiai{
\begin{itemchoice}
\item Đúng. $6\text{ m}^3 = 6\,000\text{ lít}$. Tỉ số phần trăm: $2\,400 : 6\,000 = 40\%$.
\item Đúng. Buổi chiều dùng: $2\,400 \times \dfrac{1}{3} = 800\text{ (lít)}$.
\item Sai. Lượng nước còn lại là $2\,400 - 800 = 1\,600\text{ (lít)}$.
\item Đúng. Lượng nước cần bơm thêm để đầy bể là $6\,000 - 1\,600 = 4\,400\text{ (lít)}$.
\end{itemchoice}
}
\end{ex}

\textbf{Phần 3. Câu trả lời ngắn}
\setcounterpageref{ex}{0}

\begin{ex} %[5H8V48]
Một bể hình hộp chữ nhật có diện tích đáy $12\text{ m}^2$ đang chứa $7\,200\text{ lít}$ nước. Hỏi mực nước trong bể hiện tại cao bao nhiêu đề-xi-mét?
\shortans{6}
\loigiai{
Đổi: $12\text{ m}^2 = 1\,200\text{ dm}^2$.\\
Đổi: $7\,200\text{ lít} = 7\,200\text{ dm}^3$.\\
Chiều cao của mực nước là:
$$7\,200 : 1\,200 = 6\text{ (dm)}.$$
}
\end{ex}

\begin{ex} %[5H8V48]
Một khối kim loại hình lập phương có cạnh $0{,}4\text{ m}$. Biết $1\text{ dm}^3$ kim loại đó nặng $7{,}5\text{ kg}$. Khối kim loại đó nặng bao nhiêu ki-lô-gam?
\shortans{480}
\loigiai{
Đổi $0{,}4\text{ m} = 4\text{ dm}$.\\
Thể tích khối kim loại là:
$$4 \times 4 \times 4 = 64\text{ (dm}^3\text{)}.$$
Khối kim loại đó cân nặng là:
$$64 \times 7{,}5 = 480\text{ (kg)}.$$
}
\end{ex}

\begin{ex} %[5H8V48]
Một hồ bơi có chiều dài $25\text{ m}$, chiều rộng $10\text{ m}$ và mực nước sâu trung bình $1{,}6\text{ m}$. Hỏi hồ bơi đó đang chứa bao nhiêu mét khối nước?
\shortans{400}
\loigiai{
Thể tích nước trong hồ bơi là:
$$25 \times 10 \times 1{,}6 = 400\text{ (m}^3\text{)}.$$
}
\end{ex}

\begin{ex} %[5H8V48]
Có hai bể nước hình hộp chữ nhật. Bể thứ nhất chứa được $3\,600\text{ lít}$ nước. Bể thứ hai có dung tích gấp $1{,}5$ lần bể thứ nhất. Hỏi cả hai bể chứa được tất cả bao nhiêu đề-xi-mét khối nước?
\shortans{9000}
\loigiai{
Bể thứ nhất chứa được $3\,600\text{ dm}^3$.\\
Bể thứ hai chứa được:
$$3\,600 \times 1{,}5 = 5\,400\text{ (dm}^3\text{)}.$$
Cả hai bể chứa được tất cả:
$$3\,600 + 5\,400 = 9\,000\text{ (dm}^3\text{)}.$$
}
\end{ex}

\begin{ex} %[5H8V48]
Một téc nước hình trụ chứa $2\,000\text{ lít}$ nước. Mỗi ngày gia đình dùng hết $250\text{ dm}^3$ nước. Hỏi lượng nước đó đủ cho gia đình sử dụng trong bao nhiêu ngày?
\shortans{8}
\loigiai{
Đổi $2\,000\text{ lít} = 2\,000\text{ dm}^3$.\\
Số ngày dùng hết lượng nước đó là:
$$2\,000 : 250 = 8\text{ (ngày)}.$$
}
\end{ex}

\begin{ex} %[5H8V48]
Người ta xếp $27$ khối gỗ lập phương nhỏ cạnh $10\text{ cm}$ thành một khối lập phương lớn. Thể tích của khối lập phương lớn đó là bao nhiêu đề-xi-mét khối?
\shortans{27}
\loigiai{
Mỗi khối gỗ lập phương nhỏ có cạnh $10\text{ cm} = 1\text{ dm}$.\\
Thể tích mỗi khối lập phương nhỏ là:
$$1 \times 1 \times 1 = 1\text{ (dm}^3\text{)}.$$
Thể tích khối lập phương lớn là:
$$1 \times 27 = 27\text{ (dm}^3\text{)}.$$
}
\end{ex}
